\section{Introduction}

SNARK-based signature aggregation refers to proving knowledge of the correct 
verification of many signatures in a single SNARK proof. Given the quantum threat, 
it is desirable to update the massively deployed elliptic curve based cryptosystems to 
quantum resistant alternatives. Despite post-quantum solutions to signatures and SNARKs 
both being available, there is still debate as to how best to scalably aggregate signatures 
in Ethereum. 

Falcon has emerged as a promising post-quantum alternative to existing elliptic curve 
signatures. Falcon is lattice-based and prized for its compactness and speed. Though, 
its verification algorithm introduces multiple challenges regarding the establishment of 
a computational representation amenable to arithmetization for SNARKs. Firstly, Falcon 
signature arithmetic relies upon polynomial rings with coefficients modulo $q=12289$. The 
field $\mathbb{F}_q$ is significantly smaller than those typical in modern hash-based 
SNARKs (e.g. KoalaBear). Encoding $\mathbb{F}_q$ arithmetic within the KoalaBear prime 
field requires non-native arithmetic, a process in which foreign field arithmetic is 
simulated in the native field. Doing so adds extra constraints for every foreign field 
operation compared to operations over the native field. As well, the size of $\mathbb{F}_q$ is 
less than $14$ bits, too small to permit Reed-Solomon encodings of the moderate witness 
size $2^{14}$. As a result, encoding and commitment would need to be done over an extension field, 
a performance degradation for polynomial commitments based on FRI. Furthermore, Falcon signature 
verification also employs arithmetic over the integers for $\ell_2$ norm checks.  
As a result, the SNARK must be capable of handling both integer and cyclotomic ring arithmetic 
in addition to any non-native arithmetic.

In this work we present KoalaFalcon, a variant of Falcon modified to work over the 
KoalaBear prime modulus. Our primary modifications are to fix modulus $q$ to the 
KoalaBear prime, and swap the hash function to Poseidon. This approach enables the 
ring arithmetic component of KoalaFalcon signature verification to be arithmetized at the 
coefficient level without non-native arithmetic for SNARKs defined over the KoalaBear 
prime field. As a result, SNARK proofs of aggregate KoalaFalcon signatures may rely on much 
of the same cryptographic infrastructure as existing hash-based SNARKs (e.g. LeanVM), 
making them more easily incorporated into recursive hash-based SNARKs.

We analyze how our modifications to Falcon affect the security of the resulting signature 
scheme, arguing that our changes are amenable to the same style of argument employed by 
CoreFalcon+, and retain similar security to the original Falcon scheme. The reasoning behind 
this position argues that changing the modulus in CoreFalcon+ does not impose new consequences to 
the applicability of the Rényi divergence based approach to the security reduction from CoreFalcon+ 
forgery to $t\text{-}\mathcal{R}\text{-}\mathbf{ISIS}$ under the standard CoreFalcon+ approach to 
parameterization. We define two parameter sets for KoalaFalcon achieving UF-CMA bit security estimates 
of 117.98 and 256 respectively using these techniques. As well, we observe that our estimates remain 
similar to the original Falcon scheme despite these changes, arguing why by analyzing the impact of the 
parameter change on lattice reduction attacks.

\subsection{Related Work}

\paragraph{SNARKs} 

SNARK proof systems prove and verify computational claims by combining an interactive oracle proof 
(IOP) with a polynomial commitment scheme (PCS). They are made non-interactive by simulating 
verifier responses with the Fiat-Shamir heuristic. In general, the design choices behind a 
SNARK proof system can be roughly divided into those about \emph{arithmetization} and those 
about the \emph{PCS}. Considerations of arithmetization revolve around how applications are 
compiled into arithmetic circuits, and what IOPs are used to verify that a witness attests to the 
correct execution of a computation. The PCS necessarily corresponds to how 
arithmetization is performed, as the prover must make a binding commitment to the witness of 
an arithmetized computation prior to any interaction. The commitment is later opened at a point 
determined over the course of the IOP. The field of arithmetization must be compatible enough 
with the PCS that this opening can be made in order for the soundness of the IOP to be established 
with respect to the initial commitment. Non-native arithmetic is commonly used when applications are 
defined over fields incompatible with the chosen field of arithmetization (a process where 
arithmetic for a foreign field is simulated within the field of arithmetization).

\paragraph{Arithmetic Hashes and Recursive SNARKs}

Arithmetic hashes are hash functions designed specifically to be efficiently incorporated
into arithmetic circuits. The best example being Poseidon~\cite{cryptoeprint:2019/458}, whose 
permutation function consists entirely of arithmetic operations. Poseidon has become a standard 
primitive in the construction of Recursive SNARKs, which refer to SNARK proof systems that partition 
the full execution trace of computations too large to feasibly prove within a single proof,
recursively proving those partitions. Proof recursion is achieved by starting with a proof 
of the first partition, then in each recursive proof simultaneously proves the next partition 
while verifying the previous in the same circuit.

\paragraph{Primitive-Friendly SNARKs vs SNARK-Friendly Primitives}\label{friendliness} 

SNARKs which prove the execution of cryptographic primitives that are not easily arithmetized 
often utilize \emph{SNARK-Friendly} primitives, such as Poseidon. An alternative approach to the 
non-native arithmetic problem is to use a \emph{Primitive-Friendly} SNARK. Examples include
Binius~\cite{cryptoeprint:2023/1784},
Flock~\cite{cryptoeprint:2026/1329}, and Labradoor~\cite{cryptoeprint:2022/1341}.
Such schemes define arithmetization schemes that efficiently encode applications with
special forms of arithmetic. Binius and Flock, for example, are especially efficient 
at proving the execution of 
binary circuits. 

A downside of primitive-friendly SNARKs is their bespoke nature. When the proof system is designed 
around a single arithmetization objective, it may become difficult to compose such proofs into larger 
proof systems (e.g. recursive SNARKs) without the reintroduction of non-native arithmetic. For example, 
if a primitive-friendly SNARK uses arithmetic incompatible with that of a SNARK-friendly hash,  
it is not straightforward how to combine them into a unified protocol. As well, if the PCS of the 
primitive-friendly SNARK uses arithmetic incompatible with that of the broader system, the task of 
batching PCS proofs is also no longer straightforward. Such incompatibilities introduce additional 
design and security considerations beyond those standard to composable proofs.

\paragraph{SNARK-Based XMSS Signature Aggregation}

Another approach to post-quantum signature aggregation is SNARK-based aggregation of XMSS 
signatures ~\cite{cryptoeprint:2011/484}. The XMSS approach to signatures involves taking a 
one-time signature and making it into a many-time signature by creating a Merkle root of 
their public keys. An inclusion proof alongside a one-time signature 
becomes a many-time signature, with aggregate verification provable with hash-based SNARKs 
~\cite{10.62056/aey7qjp10}. SNARK-based aggregation of XMSS signatures suffer from large circuits 
dominated by hashing. They are also stateful, requiring more considerations placed on the signer 
to remain secure in practice. 

\paragraph{Falcon Signature Aggregation with Labradoor}

Labradoor~\cite{cryptoeprint:2022/1341} is a lattice-based SNARK for proving R1CS. 
Recently, it has been used to prove Falcon signatures in aggregate 
~\cite{cryptoeprint:2024/311} and serves as the most direct performance comparison to the 
approach taken in this work.

\subsection{Contributions}

\begin{itemize}
  \item A variant of CoreFalcon+ (KoalaFalcon) over the KoalaBear 
  prime modulus using Poseidon for message hashes. We demonstrate the existence 
  of parameters that retain provable security under the Rényi 
  divergence approach to provable security. These parameters result 
  in favorable hardness estimates via the core-SVP methodology, achieving a UF-CMA bit security of $117.98$ and 
  $256$ bits for $\lambda = 128$ and $\lambda = 256$ respectively.
  We argue that aside from two conditions on the algebraic structure implied 
  by concrete parameter choices, and an adherence to a lower bound on the 
  standard deviation for signature sampling, the reduction applies for 
  general moduli.

\end{itemize}
